%%%%%%%%%%%%%%%%%%%%%%%%%%%%%%%%%%%%%%%%%
%  My documentation report
%  Objetive: Explain what I did and how, so someone can continue with the investigation
%
% Important note:
% Chapter heading images should have a 2:1 width:height ratio,
% e.g. 920px width and 460px height.
%
%%%%%%%%%%%%%%%%%%%%%%%%%%%%%%%%%%%%%%%%%


%----------------------------------------------------------------------------------------
%	PACKAGES AND OTHER DOCUMENT CONFIGURATIONS
%----------------------------------------------------------------------------------------

\documentclass[12pt]{article} % Default font size and left-justified equations

\usepackage[top=3cm,bottom=3cm,left=3.2cm,right=3.2cm,headsep=10pt,letterpaper]{geometry} % Page margins

\usepackage{xcolor} % Required for specifying colors by name
\definecolor{ocre}{RGB}{52,177,201} % Define the orange color used for highlighting throughout the book

% Font Settings
\usepackage{avant} % Use the Avantgarde font for headings
%\usepackage{times} % Use the Times font for headings
\usepackage{mathptmx} % Use the Adobe Times Roman as the default text font together with math symbols from the Sym­bol, Chancery and Com­puter Modern fonts
\usepackage{microtype} % Slightly tweak font spacing for aesthetics
\usepackage[utf8]{inputenc} % Required for including letters with accents
\usepackage[T1]{fontenc} % Use 8-bit encoding that has 256 glyphs

% Bibliography
\usepackage[style=apa,sorting=nyt,sortcites=true,autopunct=true,babel=hyphen,hyperref=true,abbreviate=false,backref=true,backend=biber]{biblatex}

\addbibresource{bibliography.bib} % BibTeX bibliography file
\defbibheading{bibempty}{}

%----------------------------------------------------------------------------------------
%	VARIOUS REQUIRED PACKAGES
%----------------------------------------------------------------------------------------

\usepackage{titlesec} % Allows customization of titles

\usepackage{graphicx} % Required for including pictures
\graphicspath{{Pictures/}} % Specifies the directory where pictures are stored
% \graphicspath{{Plots/}}
\usepackage{lipsum} % Inserts dummy text

\usepackage{tikz} % Required for drawing custom shapes

\usepackage[english]{babel} % English language/hyphenation

\usepackage{enumitem} % Customize lists
\setlist{nolistsep} % Reduce spacing between bullet points and numbered lists

\usepackage{booktabs} % Required for nicer horizontal rules in tables

\usepackage{eso-pic} % Required for specifying an image background in the title page

%----------------------------------------------------------------------------------------
%	MAIN TABLE OF CONTENTS
%----------------------------------------------------------------------------------------

\usepackage{titletoc} % Required for manipulating the table of contents

\contentsmargin{0cm} % Removes the default margin
% Chapter text styling
\titlecontents{chapter}[1.25cm] % Indentation
{\addvspace{15pt}\large\sffamily\bfseries} % Spacing and font options for chapters
{\color{ocre!60}\contentslabel[\Large\thecontentslabel]{1.25cm}\color{ocre}} % Chapter number
{}  
{\color{ocre!60}\normalsize\sffamily\bfseries\;\titlerule*[.5pc]{.}\;\thecontentspage} % Page number
% Section text styling
\titlecontents{section}[1.25cm] % Indentation
{\addvspace{5pt}\sffamily\bfseries} % Spacing and font options for sections
{\contentslabel[\thecontentslabel]{1.25cm}} % Section number
{}
{\sffamily\hfill\color{black}\thecontentspage} % Page number
[]
% Subsection text styling
\titlecontents{subsection}[1.25cm] % Indentation
{\addvspace{1pt}\sffamily\small} % Spacing and font options for subsections
{\contentslabel[\thecontentslabel]{1.25cm}} % Subsection number
{}
{\sffamily\;\titlerule*[.5pc]{.}\;\thecontentspage} % Page number
[] 

%----------------------------------------------------------------------------------------
%	MINI TABLE OF CONTENTS IN CHAPTER HEADS
%----------------------------------------------------------------------------------------

% Section text styling
\titlecontents{lsection}[0em] % Indendating
{\footnotesize\sffamily} % Font settings
{}
{}
{}

% Subsection text styling
\titlecontents{lsubsection}[.5em] % Indentation
{\normalfont\footnotesize\sffamily} % Font settings
{}
{}
{}
 
%----------------------------------------------------------------------------------------
%	PAGE HEADERS
%----------------------------------------------------------------------------------------

\usepackage{fancyhdr} % Required for header and footer configuration

\pagestyle{fancy}
%\renewcommand{\chaptermark}[1]{\markboth{\sffamily\normalsize\bfseries\chaptername\ \thechapter.\ #1}{}} % Chapter text font settings

\renewcommand{\sectionmark}[1]{\markright{\sffamily\normalsize\thesection\hspace{5pt}#1}{}} % Section text font settings
\fancyhf{} \fancyhead[LE,RO]{\sffamily\normalsize\thepage} % Font setting for the page number in the header
\fancyhead[LO]{\rightmark} % Print the nearest section name on the left side of odd pages
\fancyhead[RE]{\leftmark} % Print the current chapter name on the right side of even pages
\renewcommand{\headrulewidth}{0.5pt} % Width of the rule under the header
\addtolength{\headheight}{2.5pt} % Increase the spacing around the header slightly
\renewcommand{\footrulewidth}{0pt} % Removes the rule in the footer
\fancypagestyle{plain}{\fancyhead{}\renewcommand{\headrulewidth}{0pt}} % Style for when a plain pagestyle is specified

% Removes the header from odd empty pages at the end of chapters
\makeatletter
\renewcommand{\cleardoublepage}{
\clearpage\ifodd\c@page\else
\hbox{}
\vspace*{\fill}
\thispagestyle{empty}
%\newpage
\fi}

%----------------------------------------------------------------------------------------
%	THEOREM STYLES
%----------------------------------------------------------------------------------------

\usepackage{amsmath,amsfonts,amssymb,amsthm} % For math equations, theorems, symbols, etc

\newcommand{\intoo}[2]{\mathopen{]}#1\,;#2\mathclose{[}}
\newcommand{\ud}{\mathop{\mathrm{{}d}}\mathopen{}}
\newcommand{\intff}[2]{\mathopen{[}#1\,;#2\mathclose{]}}
%\newtheorem{notation}{Notation}[chapter]

%%%%%%%%%%%%%%%%%%%%%%%%%%%%%%%%%%%%%%%%%%%%%%%%%%%%%%%%%%%%%%%%%%%%%%%%%%%
%%%%%%%%%%%%%%%%%%%% dedicated to boxed/framed environements %%%%%%%%%%%%%%
%%%%%%%%%%%%%%%%%%%%%%%%%%%%%%%%%%%%%%%%%%%%%%%%%%%%%%%%%%%%%%%%%%%%%%%%%%%
\newtheoremstyle{ocrenumbox}% % Theorem style name
{0pt}% Space above
{0pt}% Space below
{\normalfont}% % Body font
{}% Indent amount
{\small\bf\sffamily\color{ocre}}% % Theorem head font
{\;}% Punctuation after theorem head
{0.25em}% Space after theorem head
{\small\sffamily\color{ocre}\thmname{#1}\nobreakspace\thmnumber{\@ifnotempty{#1}{}\@upn{#2}}% Theorem text (e.g. Theorem 2.1)
\thmnote{\nobreakspace\the\thm@notefont\sffamily\bfseries\color{black}---\nobreakspace#3.}} % Optional theorem note
\renewcommand{\qedsymbol}{$\blacksquare$}% Optional qed square

\newtheoremstyle{blacknumex}% Theorem style name
{5pt}% Space above
{5pt}% Space below
{\normalfont}% Body font
{} % Indent amount
{\small\bf\sffamily}% Theorem head font
{\;}% Punctuation after theorem head
{0.25em}% Space after theorem head
{\small\sffamily{\tiny\ensuremath{\blacksquare}}\nobreakspace\thmname{#1}\nobreakspace\thmnumber{\@ifnotempty{#1}{}\@upn{#2}}% Theorem text (e.g. Theorem 2.1)
\thmnote{\nobreakspace\the\thm@notefont\sffamily\bfseries---\nobreakspace#3.}}% Optional theorem note

\newtheoremstyle{blacknumbox} % Theorem style name
{0pt}% Space above
{0pt}% Space below
{\normalfont}% Body font
{}% Indent amount
{\small\bf\sffamily}% Theorem head font
{\;}% Punctuation after theorem head
{0.25em}% Space after theorem head
{\small\sffamily\thmname{#1}\nobreakspace\thmnumber{\@ifnotempty{#1}{}\@upn{#2}}% Theorem text (e.g. Theorem 2.1)
\thmnote{\nobreakspace\the\thm@notefont\sffamily\bfseries---\nobreakspace#3.}}% Optional theorem note

%%%%%%%%%%%%%%%%%%%%%%%%%%%%%%%%%%%%%%%%%%%%%%%%%%%%%%%%%%%%%%%%%%%%%%%%%%%
%%%%%%%%%%%%% dedicated to non-boxed/non-framed environements %%%%%%%%%%%%%
%%%%%%%%%%%%%%%%%%%%%%%%%%%%%%%%%%%%%%%%%%%%%%%%%%%%%%%%%%%%%%%%%%%%%%%%%%%
\newtheoremstyle{ocrenum}% % Theorem style name
{5pt}% Space above
{5pt}% Space below
{\normalfont}% % Body font
{}% Indent amount
{\small\bf\sffamily\color{ocre}}% % Theorem head font
{\;}% Punctuation after theorem head
{0.25em}% Space after theorem head
{\small\sffamily\color{ocre}\thmname{#1}\nobreakspace\thmnumber{\@ifnotempty{#1}{}\@upn{#2}}% Theorem text (e.g. Theorem 2.1)
\thmnote{\nobreakspace\the\thm@notefont\sffamily\bfseries\color{black}---\nobreakspace#3.}} % Optional theorem note
\renewcommand{\qedsymbol}{$\blacksquare$}% Optional qed square
\makeatother

% Defines the theorem text style for each type of theorem to one of the three styles above
\newcounter{dummy} 
\numberwithin{dummy}{section}
\theoremstyle{ocrenumbox}


\newtheorem{theoremeT}[dummy]{Theorem}
\newtheorem{lemma}[dummy]{Lemma}
\newtheorem{observation}[dummy]{Observation}
\newtheorem{proposition}[dummy]{Proposition}
% \newtheorem{definition}[dummy]{Definition}
\newtheorem{claim}[dummy]{Claim}
\newtheorem{fact}[dummy]{Fact}
\newtheorem{assumption}[dummy]{Assumption}

%\newtheorem{problem}{Problem}[chapter]
% \newtheorem{exercise}{Exercise}[chapter]
\theoremstyle{blacknumex}
\newtheorem{exampleT}{Example}[section]
\theoremstyle{blacknumbox}
%\newtheorem{vocabulary}{Vocabulary}[chapter]
\newtheorem{definitionT}{Definition}[section]
\newtheorem{corollaryT}[dummy]{Corollary}
\theoremstyle{ocrenum}

%----------------------------------------------------------------------------------------
%	DEFINITION OF COLORED BOXES
%----------------------------------------------------------------------------------------

\RequirePackage[framemethod=default]{mdframed} % Required for creating the theorem, definition, exercise and corollary boxes

% Theorem box
\newmdenv[skipabove=7pt,
skipbelow=7pt,
backgroundcolor=black!5,
linecolor=ocre,
innerleftmargin=5pt,
innerrightmargin=5pt,
innertopmargin=5pt,
leftmargin=0cm,
rightmargin=0cm,
innerbottommargin=5pt]{tBox}

% Exercise box	  
\newmdenv[skipabove=7pt,
skipbelow=7pt,
rightline=false,
leftline=true,
topline=false,
bottomline=false,
backgroundcolor=ocre!10,
linecolor=ocre,
innerleftmargin=5pt,
innerrightmargin=5pt,
innertopmargin=5pt,
innerbottommargin=5pt,
leftmargin=0cm,
rightmargin=0cm,
linewidth=4pt]{eBox}	

% Definition box
\newmdenv[skipabove=7pt,
skipbelow=7pt,
rightline=false,
leftline=true,
topline=false,
bottomline=false,
linecolor=ocre,
innerleftmargin=5pt,
innerrightmargin=5pt,
innertopmargin=0pt,
leftmargin=0cm,
rightmargin=0cm,
linewidth=4pt,
innerbottommargin=0pt]{dBox}	

% Corollary box
\newmdenv[skipabove=7pt,
skipbelow=7pt,
rightline=false,
leftline=true,
topline=false,
bottomline=false,
linecolor=gray,
backgroundcolor=black!5,
innerleftmargin=5pt,
innerrightmargin=5pt,
innertopmargin=5pt,
leftmargin=0cm,
rightmargin=0cm,
linewidth=4pt,
innerbottommargin=5pt]{cBox}

% Creates an environment for each type of theorem and assigns it a theorem text style from the "Theorem Styles" section above and a colored box from above
\newenvironment{theorem}{\begin{tBox}\begin{theoremeT}}{\end{theoremeT}\end{tBox}}
\newenvironment{exercise}{\begin{eBox}\begin{exerciseT}}{\hfill{\color{ocre}\tiny\ensuremath{\blacksquare}}\end{exerciseT}\end{eBox}}				  
\newenvironment{definition}{\begin{dBox}\begin{definitionT}}{\end{definitionT}\end{dBox}}	
\newenvironment{example}{\begin{exampleT}}{\hfill{\tiny\ensuremath{\blacksquare}}\end{exampleT}}		
\newenvironment{corollary}{\begin{cBox}\begin{corollaryT}}{\end{corollaryT}\end{cBox}}	

%----------------------------------------------------------------------------------------
%	REMARK ENVIRONMENT
%----------------------------------------------------------------------------------------

\newenvironment{remark}{\par\vspace{10pt}\small % Vertical white space above the remark and smaller font size
\begin{list}{}{
\leftmargin=35pt % Indentation on the left
\rightmargin=25pt}\item\ignorespaces % Indentation on the right
\makebox[-2.5pt]{\begin{tikzpicture}[overlay]
\node[draw=ocre!60,line width=1pt,circle,fill=ocre!25,font=\sffamily\bfseries,inner sep=2pt,outer sep=0pt] at (-15pt,0pt){\textcolor{ocre}{R}};\end{tikzpicture}} % Orange R in a circle
\advance\baselineskip -1pt}{\end{list}\vskip5pt} % Tighter line spacing and white space after remark

%----------------------------------------------------------------------------------------
%	SECTION NUMBERING IN THE MARGIN
%----------------------------------------------------------------------------------------

\makeatletter
\renewcommand{\@seccntformat}[1]{\llap{\textcolor{ocre}{\csname the#1\endcsname}\hspace{1em}}}                    
\renewcommand{\section}{\@startsection{section}{1}{\z@}
{-4ex \@plus -1ex \@minus -.4ex}
{1ex \@plus.2ex }
{\normalfont\large\sffamily\bfseries}}
\renewcommand{\subsection}{\@startsection {subsection}{2}{\z@}
{-3ex \@plus -0.1ex \@minus -.4ex}
{0.5ex \@plus.2ex }
{\normalfont\sffamily\bfseries}}
\renewcommand{\subsubsection}{\@startsection {subsubsection}{3}{\z@}
{-2ex \@plus -0.1ex \@minus -.2ex}
{.2ex \@plus.2ex }
{\normalfont\small\sffamily\bfseries}}                        
\renewcommand\paragraph{\@startsection{paragraph}{4}{\z@}
{-2ex \@plus-.2ex \@minus .2ex}
{.1ex}
{\normalfont\small\sffamily\bfseries}}

%----------------------------------------------------------------------------------------
%	HYPERLINKS IN THE DOCUMENTS
%----------------------------------------------------------------------------------------

% For an unclear reason, the package should be loaded now and not later
\usepackage{hyperref}
\hypersetup{hidelinks,backref=true,pagebackref=true,hyperindex=true,colorlinks=false,breaklinks=true,urlcolor= ocre,bookmarks=true,bookmarksopen=false,pdftitle={Title},pdfauthor={Author}}

%----------------------------------------------------------------------------------------
%	CHAPTER HEADINGS
%----------------------------------------------------------------------------------------

% The set-up below should be (sadly) manually adapted to the overall margin page septup controlled by the geometry package loaded in the main.tex document. It is possible to implement below the dimensions used in the goemetry package (top,bottom,left,right)... TO BE DONE

\newcommand{\thechapterimage}{}
\newcommand{\chapterimage}[1]{\renewcommand{\thechapterimage}{#1}}

% Numbered chapters with mini tableofcontents
\def\thechapter{\arabic{chapter}}
\def\@makechapterhead#1{
\thispagestyle{empty}
{\centering \normalfont\sffamily
\ifnum \c@secnumdepth >\m@ne
\if@mainmatter
\startcontents
\begin{tikzpicture}[remember picture,overlay]
\node at (current page.north west)
{\begin{tikzpicture}[remember picture,overlay]
\node[anchor=north west,inner sep=0pt] at (0,0) {\includegraphics[width=\paperwidth]{\thechapterimage}};
%%%%%%%%%%%%%%%%%%%%%%%%%%%%%%%%%%%%%%%%%%%%%%%%%%%%%%%%%%%%%%%%%%%%%%%%%%%%%%%%%%%%%
% Commenting the 3 lines below removes the small contents box in the chapter heading
%\fill[color=ocre!10!white,opacity=.6] (1cm,0) rectangle (8cm,-7cm);
%\node[anchor=north west] at (1.1cm,.35cm) {\parbox[t][8cm][t]{6.5cm}{\huge\bfseries\flushleft \printcontents{l}{1}{\setcounter{tocdepth}{2}}}};
\draw[anchor=west] (5cm,-9cm) node [rounded corners=20pt,fill=ocre!10!white,text opacity=1,draw=ocre,draw opacity=1,line width=1.5pt,fill opacity=.6,inner sep=12pt]{\huge\sffamily\bfseries\textcolor{black}{\thechapter. #1\strut\makebox[22cm]{}}};
%%%%%%%%%%%%%%%%%%%%%%%%%%%%%%%%%%%%%%%%%%%%%%%%%%%%%%%%%%%%%%%%%%%%%%%%%%%%%%%%%%%%%
\end{tikzpicture}};
\end{tikzpicture}}
\par\vspace*{230\p@}
\fi
\fi}

% Unnumbered chapters without mini tableofcontents (could be added though) 
\def\@makeschapterhead#1{
\thispagestyle{empty}
{\centering \normalfont\sffamily
\ifnum \c@secnumdepth >\m@ne
\if@mainmatter
\begin{tikzpicture}[remember picture,overlay]
\node at (current page.north west)
{\begin{tikzpicture}[remember picture,overlay]
\node[anchor=north west,inner sep=0pt] at (0,0) {\includegraphics[width=\paperwidth]{\thechapterimage}};
\draw[anchor=west] (5cm,-9cm) node [rounded corners=20pt,fill=ocre!10!white,fill opacity=.6,inner sep=12pt,text opacity=1,draw=ocre,draw opacity=1,line width=1.5pt]{\huge\sffamily\bfseries\textcolor{black}{#1\strut\makebox[22cm]{}}};
\end{tikzpicture}};
\end{tikzpicture}}
\par\vspace*{230\p@}
\fi
\fi
}
\makeatother % Insert the commands.tex file which contains the majority of the structure behind the template

%----------------------------------------------------------------------------------------
%	Definitions of new commands
%----------------------------------------------------------------------------------------

\def\R{\mathbb{R}}
\newcommand{\cvx}{convex}
\begin{document}

%----------------------------------------------------------------------------------------
%	TITLE PAGE
%----------------------------------------------------------------------------------------

\begingroup
\thispagestyle{empty}
\AddToShipoutPicture*{\put(0,0){\includegraphics[scale=1.05]{Pictures/covveeeer12.PNG}}} % Image background
\centering{
\vspace*{2.5cm}
\par\normalfont\fontsize{35}{35}\sffamily\selectfont
\textbf{Lecture Notes on Separation Axioms}\\
{\LARGE Math 6201 - Topology } \\ % Book title
\vspace{-0.3cm}{\Large AY 2024-2025 Term 1} \\
{\Large  James Israel B. Montillano} \par \vspace{0.2 cm}  % Author name
{\Large Romelyn M. Dumanhog} } \\~\\
{\normalsize Department of Computer, Information Sciences, and Mathematics} \\
{\small (May 2025)} \\
\endgroup

%----------------------------------------------------------------------------------------
%	COPYRIGHT PAGE
%----------------------------------------------------------------------------------------

%\newpage
%~\vfill
%\thispagestyle{empty}

%%\noindent Copyright \copyright\ 2014 Andrea Hidalgo\\ % Copyright notice

%\noindent \textsc{Summer Research Internship, University of Western Ontario}\\

%\noindent \textsc{github.com/LaurethTeX/Clustering}\\ % URL

%\noindent This research was done under the supervision of Dr. Pauline Barmby with the financial support of the MITACS Globalink Research Internship Award within a total of 12 weeks, from June 16th to September 5th of 2014.\\ % License information

%\noindent \textit{First release, August 2014} % Printing/edition date

%----------------------------------------------------------------------------------------
%	TABLE OF CONTENTS
%----------------------------------------------------------------------------------------

\chapterimage{head1.png} % Table of contents heading image

\pagestyle{empty} % No headers
\newpage
\tableofcontents % Print the table of contents itself

%\cleardoublepage % Forces the first chapter to start on an odd page so it's on the right

\pagestyle{fancy} % Print headers again

%----------------------------------------------------------------------------------------
%	CHAPTER 1
%----------------------------------------------------------------------------------------

\newpage
%I edit nya ni romelyn kay kang cherlyn ni.
\section{Preliminaries}
This section reviews some fundamental concepts needed to understand the basic parts on separation axioms on topology. 

\subsection{Sets}
\begin{definition}[Subset]\parencite{lipschutz1965} A set $A$ is a subset of a set $B$ or, equivalently, $B$ is a superset of $A$, written $A\subset B$ or $B\supset A$ iff each element in $A$ also belongs to $B$; that is, if $x\in A$ implies $x\in B$.
\end{definition}

\begin{definition}[Union]\parencite{lipschutz1965}
The union of two sets $A$ and $B$, denoted by $A\cup B$, is the set of all elements which belong to $A$ or $B$, i.e., $A\cup B = \{x:x\in A \quad \text{or} \quad x\in B\}$     
\end{definition}

\begin{definition}[Intersection]\parencite{lipschutz1965}
The intersection of two sets $A$ and $B$, denoted by $A\cap B$, is the set of elements which belong to both $A$ and $B$, i.e., $A\cap B=\{x:x\in A \quad \text{and} \quad x\in B)$   
\end{definition}


\begin{definition}[Power Sets]\parencite{Dugundji1966}
Let $A$ be any set. Its power set $\mathcal{P}(A)$ is the set of all subsets of $A$.
\end{definition}

\begin{definition}[Countable]\parencite{Dugundji1966}
 A set $A$ is countable if it is finite or equivalent to the set $\mathbb{N}$ of counting
 numbers. If $A\equiv N$, then $A$ is called countably infinite or denumerable.
\end{definition}


\subsection{Functions, or Maps}

\begin{definition}[Map] \parencite{Dugundji1966}
Let $X$ and $Y$ be two sets. A \textit{map} $f: X \rightarrow Y$ \textit{(or function with domain $X$ and range $Y$)} is a subset $f\subset X \times Y$ with the property: for each $x\in X$, there is one, and only one, $y\in Y$ satisfying $(x,y)\in f$.To denote $(x,y)\in f$, we write $y=f(x)$ and say that $y$ is the image of $x$ under $f$.
\end{definition}



\begin{definition}[Surjective function]\parencite{holmes_topology}
Let $f : X \to Y$ be a function. We say that $f$ is surjective if for each $y \in Y$ , there exists $x \in X$ such that $f(x) = y$
\end{definition}

\begin{definition}[Injective function]\parencite{holmes_topology}
Let $f : X \to Y$ be a function. We say that
$f$ is injective if $f(x_1) = f(x_2)$ implies $x_1 = x_2 (x_i \in X)$
\end{definition}


\begin{definition}[Bijective function]\parencite{holmes_topology}
Let \( f : X \to Y \) be a function. We say that \( f \) is bijective if it is both injective and surjective.
\end{definition}

%\begin{example}
%$S_1^n  \{ X | X=X^n ,\lambda(x) \ge 0\}$\\
%A matrix with positive eigenvalues.
%\end{example}

\subsection{Topology}
\begin{definition}[Topology]\parencite{Dugundji1966} 
Let $X$ be a set. A \textit{topology} (or topological structure) in $X$ is a family $\tau $ of subsets of $X$ that satisfies:
\begin{enumerate}
    \item [(1).] Each union of members of $\tau$ is also a member of $\tau$.
    \item [(2).] Each \textit{finite} intersection of members of $\tau$ is also a member of $\tau$.
    \item [(3).] $\emptyset$ and $X$ are members of $\tau$.
\end{enumerate}
\end{definition}

\begin{example}[Discrete topology]
 Let $X$ be any set and 
 $\tau=\mathcal{P}(X)$. Then $\tau$ is a topology on $X$.
 \begin{proof}  ~\\
 \begin{enumerate}
     \item Clearly, $X$, $\varnothing \in \tau$
     \item Since every possible subset of $X$ is included in $\tau$, the union of any combination of these subsets will always result in another subset of $X$, which is already in $\tau$.
     \item Since all subsets are open, the intersection of any finite number of them will still be a subset of $X$, hence in $\tau$.
 \end{enumerate}
 Thus, $\tau$ is a topology on $X$.
 \end{proof}
\end{example}

\begin{example}[Indiscrete topology]
Let $X$ be any set and $\tau=\{X,\varnothing\}$.  Then $\tau$ is a topology on $X$
\begin{proof}  ~\\
\begin{enumerate}
    \item $X \in \tau$ and $\varnothing \in \tau$
    \item $X \cup X=X \cup \varnothing= X \in \tau$\\
     $\varnothing \cup \varnothing = \varnothing \in \tau$
    \item $X \cap X= X \in \tau$\\
    $X \cap \varnothing=\varnothing \cap \varnothing = \varnothing \in \tau$
\end{enumerate}
Thus, $\tau$ is a topology on $X$.
\end{proof}
\end{example}


\begin{definition}[Open sets]\parencite{morris2020topology} 
Let $(X,\tau)$ be a topological space. Then the members of $\tau$ are said to be \emph{open sets}.
\end{definition}

\begin{definition}[Closed sets]\parencite{morris2020topology}
Let $(X,\tau)$ be a topological space. A subset $S$ of $X$ is said to be \emph{closed set} in $(X,\tau)$ if its complement in $X$, namely $X \setminus S$, is open in $(X,\tau)$.
\end{definition}

\begin{definition}[Neighborhood]\parencite{holmes_topology}
Let $x \in X$. Any open set containing $x$ is called a neighborhood of $x$.
\end{definition}

\subsection{$G_{\delta}$}
\begin{definition} \parencite{Dugundji1966}
    A set $G$ is called $G_\delta$ if it is the intersection of at most countably many open sets.
\end{definition}
\subsection{Relativization}
\begin{definition}[Subspace topology]\parencite{Dugundji1966}
 Let $(X,\tau )$ be a topological space and $Y \subset X$. The induced topology $\tau_Y$ on
 $Y$ is $\{Y\cap U : U \in \tau\}$. The pair $(Y,\tau_Y)$ is called a subspace of $(X,\tau )$.
\end{definition}

\subsection{Continuous maps and Homeomorphisms}
\begin{definition}[Continuous maps]\parencite{BellezaCM2025}
Let \( (X, \tau_X) \) and \( (Y, \tau_Y) \) be topological spaces. A map \( f : X \to Y \) is called \emph{continuous} if the inverse image of each open set in \( Y \) is open in \( X \).  
That is, \( f^{-1}\) maps $\tau_Y \to \tau_X$.   
\end{definition}

\begin{example}
 Let $(X,\tau)$ be any topological space and $f : (X,\tau) \to (X,\tau)$ is defined
 by $f(x) = x$ for all $x\in X$. Then f is continuous.
\end{example}
\begin{proof}
To show that \( f \) is continuous, we need to verify that the inverse image of every open set in  
\( Y = (X, \tau) \) is open in \( X = (X, \tau) \).

Let \( O \) be an arbitrary open set in \( Y \). By definition, $f^{-1}(O) = \{ x \in X \mid f(x) \in O \}$.
Since \( f(x) = x \), this simplifies to $f^{-1}(O) = \{ x \in X \mid x \in O \} = O$. Since \( O \) is open in \( Y \), and \( Y \) has the same topology as \( X \), \( O \) is also open in \( X \). Thus, \( f^{-1}(O) = O \) is open in \( X \). Since the inverse image of every open set in \( Y \) is open in \( X \), the identity map \( f \) is continuous.
\end{proof}

\begin{remark}
Not every identity function is continuous. To see this, let \( X = \{1, 2, 3, 4\} \), and define two topologies on \( X \), $\tau_1 = \{\emptyset, \{1\}, \{3\}, \{1,3\}, \{2,3\}, \{1,2,3\}, X\}$, \\ $\tau_2 = \{\emptyset, \{1\}, \{3\}, \{1,3\}, \{2\}, \{1,2,3\}, X\}$. Clearly, \( \tau_1 \) and \( \tau_2 \) are topologies on \( X \). Then \( i : (X, \tau_1) \to (X, \tau_2) \) defined by $i(x) = x$ for all $x \in X$ is not continuous, since there exists \( \{2\} \in \tau_2 \) such that $i^{-1}(\{2\}) = \{2\ \notin \tau_1$.
\end{remark}

\begin{definition}[Homeomorphism]\parencite{Dugundji1966}
    A continuous bijective map $f: X \to Y$, such that $f^{-1}: Y \to X$ is also continuous, is called a homeomorphism (or a bicontinuous bijection) and denoted by $f: X \cong Y$. Two spaces $X, Y$ are homeomorphic, written $X \cong Y$, if there is a homeomorphism $f: X \cong Y$.
\end{definition}

\begin{example}
The identity map $1 : X\to X$ is a homeomorphism.
\end{example}

\begin{example}
Two discrete spaces $X$ and $Y$ (similarly, for indiscrete spaces), are homeomorphic if and only if there is a one-to-one function on $X$ onto $Y$ .
\end{example}
\begin{proof}
In a discrete space, every subset is open. Let \( f : X \to Y \) be a bijection (one-to-one and onto). We will show that \( f \) is a homeomorphism. To show \( f \) is continuous, the preimage of any open set in \( Y \) must be open in \( X \). Since \( Y \) is discrete, every subset of \( Y \) is open. Thus, for any open \( V \subseteq Y \),  \( f^{-1}(V) \subseteq X \). Because \( X \) is discrete, \( f^{-1}(V) \) is open in \( X \). Hence, \( f \) is continuous. Similarly, \( f^{-1} : Y \to X \) is a bijection. For any open \( U \subseteq X \), $(f^{-1})^{-1}(U) = f(U) \subseteq Y$. Since \( X \) is discrete, \( U \) is open, and since \( Y \) is discrete, \( f(U) \) is open in \( Y \).  
Thus, \( f^{-1} \) is continuous. Hence, \( f \) is a bijection, continuous, and its inverse \( f^{-1} \) is also continuous. Therefore, \( f \) is a homeomorphism, and \( X \cong Y \).
\end{proof}

\begin{definition}[Topological Invariant]
    We call any property of spaces a topological invariant if whenever it is true for one space $X$, it is also true for every space homemorphic to $X$
\end{definition}



\section{Separation Axioms}
Separation axioms are a family of topological invariants that give us new ways of distinguishing
between various spaces. The idea is to look how open sets in a space can be used to create “buffer
zones” separating pairs of points and closed sets. Separations axioms are denoted by $T_0$ $T_1$, $T_2$, etc., where
$T$ comes from the German word Trennungsaxiom, which just means “separation axiom”  \parencite{schechter1996handbook}.
\subsection{$T_0$ and $T_1$ space}
\begin{definition}[$T_0$ space or Kolmogorov Space] \parencite{milewski1994topology} 
    The space \((X, \mathcal{T})\) is said to be a \(T_0\)-space, if for any two distinct \(a, b \in X\), there is a neighborhood of at least one, which does not contain the other 
\end{definition}

\begin{example}
    Let \( X = \{a, b, c\} \) with topology \( \tau = \{\varnothing, X, \{a\}, \{b\}, \{a, b\}\} \) defined on \( X \), 
then \( (X, \tau) \) is a \( T_0 \) space
\end{example}

\begin{proof} ~\\
\begin{enumerate}
    \item for \( a \) and \( b \), there exists an open set \( \{a\} \) such that \( a \in \{a\} \) and \( b \notin \{a\} \)
    \item for \( a \) and \( c \), there exists an open set \( \{a\} \) such that \( a \in \{a\} \) and \( c \notin \{a\} \)
    \item for \( b \) and \( c \), there exists an open set \( \{b\} \) such that \( b \in \{b\} \) and \( c \notin \{b\} \)
\end{enumerate}
\end{proof}

\begin{example}
    The space $X=\{a,b\}$ with the indiscrete topology is not a $T_0$ space
\end{example}

   \begin{proof}
        The two distinct points $a$ and $b$  in $X$ is contained only in the open set $X$. Thus, it is not a $T_0$ space.
    \end{proof}

\begin{definition}[$T_1$-space or Fréchet Space ]\parencite{milewski1994topology} A topological space $(X,\mathcal{T})$ is called a \(T_1\)-space, if every single element set is closed, that is, $\forall a\in X, \{a\}=\overline{\{a\}}$ .
\end{definition}

\begin{theorem} \label{T1th}
    A topological space $(X,\tau)$ is a $T_1$ space if and only if for any pair of distinct points $a,b \in X$, the open sets $G,H \in \tau$ exist, such that $a\in G$, $b \notin G$ and $b\in H$, $a\notin H$
\end{theorem}
\begin{proof}
    Suppose $X$ is a $T_1$-space. Then, for any $x\in X$, $\{x\}$ is a closed set. Let $a,b\in X$ and $a\neq b$. The sets $X-\{a\}$ and $X-\{b\}$ are open, and 

\[
a \in X - \{b\} \quad \text{and} \quad b \notin X - \{b\}
\]
\[
b \in X - \{a\} \quad \text{and} \quad a \notin X - \{a\}.
\]

Conversely, suppose \( x \in X \). We shall show that \( \{x\} \) is closed, i.e., \( X - \{x\} \) is open.
Let \( y \in X - \{x\} \), then \( y \ne x \) and an open set \( H_y \) exists, such that

\[
y \in H_y \quad \text{and} \quad x \notin H_y.
\]

Thus,

\[
y \in H_y \subseteq X - \{x\} \quad \text{and} \quad X - \{x\} = \bigcup_{y \ne x} H_y.
\]

Since all \( H_y \) are open sets, \( X - \{x\} \) is open and \( \{x\} \) is closed, \( \{x\} = \overline{\{x\}} \).


\end{proof}

\begin{example} Consider the set $X=\{a,b,c\}$ with the cofinite topology, $$\tau=\{X,\emptyset,{a},{b},{c},\{a,b\},\{a,c\},\{b,c\}\}.$$ Verify that $(X,\tau)$ is a $T_1$ space.
\end{example}
\begin{proof} The complement of $\{a\}$ is $\{b,c\}$, which is open (its complement $\{a\}$ is finite). Similarly, $\{b\}^c=\{a,c\}$ and $\{c\}^c=\{a,b\}$ are open. By Definition 2.2, since every singleton is closed, $X$ is $T_1$.
For $a$ and $b$, there exist open sets $\{a\}$ and $\{b\}$ such that $a \in \{a\}$, $b\notin \{a\}$ and $b \in \{b\}$, $a\notin \{b\}$. This satisfies the condition in Theorem 2.1, confirming $X$ is $T_1$.

Thus, $(X,\tau)$ is a $T_1$ space.
\end{proof}

\begin{example}
The Sierpinski Space is $T_0$ but not $T_1$
\end{example}
\begin{proof}
      Recall that a \textit{{Sierpiński space}} is the topological space \( X = \{x, y\} \) with the topology given by \( \{X, \{x\}, \emptyset\} \). It is $T_0$ because for $x$ and $y$ the open set $\{x\}$ contains $x$ but not $y$. It is not $T_1$ because every open set $U$ containing $y$ (which is only $X$) contains $x$.

\end{proof}

\begin{theorem}\parencite{milewski1994topology}
    A $T_1$ space is also a $T_0$ space. \label{t1t0}
\end{theorem}
\begin{proof}
    Let $X$ be a $T_1$ space then clearly from its definition it follows that it is also a $T_0$ space. Since with any pair $a,b \in X$ there exist an open set $G$ with $a\in G$ and $b\notin G$.
\end{proof}

\begin{theorem}\parencite{milewski1994topology} If $(X,\tau)$ and $(Y,\tau')$ are homeomorphic  and $(X,\tau)$ is a $T_1$-space (or $T_0$) then so is  $(Y,\tau')$ 
\end{theorem}
\begin{proof}
    Let $f$ denote a homeomorphism
\[ f: X \to Y \]
and $X$ be a $T_1$-space. A space $(X, \tau)$ is $T_1$, if and only if every one-point subset of $X$ is closed.
Let $y$ represent any point of $Y$, $y \in Y$. The set $f^{-1}(y)$ is a one-point subset of $X$ and since $X$ is $T_1$, the set $\{f^{-1}(y)\}$ is closed.

Since $f: X \to Y$ is a homeomorphism, it maps closed sets into closed sets.
Therefore, for any $y \in Y$
\[ \{y\} = \overline{\{y\}}. \]
Thus, $(Y, \tau)$ is a $T_1$-space. Similarl proof can be done to when if $X$ is $T_0$, then so is $Y$.
\end{proof}


\subsection{Hausdorff spaces}
\begin{definition}[$T_2$-space or Hausdorff space] \parencite{Dugundji1966} A space $X$ is Hausdorff (or separated) if each two distinct points have nonintersecting nbds, that is, whenever, $p \neq q$ there are nbds $U(p), V(q)$ such that $U \cap V = \varnothing$. 
\end{definition}

\begin{example} The real line $R$ with the standard topology is a Hausdorff space.
\end{example}
 \begin{proof}
 For any two distinct points \( x, y \in \mathbb{R} \), let \( d = |x - y| > 0 \).  
Then the open intervals  
\[ 
U = \left( x - \frac{d}{2}, \, x + \frac{d}{2} \right) \quad \text{and} \quad V = \left( y - \frac{d}{2}, \, y + \frac{d}{2} \right) 
\]  
are disjoint neighborhoods of \( x \) and \( y \), respectively.  
Thus, \( \mathbb{R} \) is Hausdorff.
    \end{proof}

\begin{theorem} \parencite{milewski1994topology}
    Each Hausdorff space is a $T_1$ space. \label{hto1}
\end{theorem}
\begin{proof}
    Suppose $X$ is a Hausdorff space. Then from the definition, for two distinct points $x$ and $y$, there exist two open sets $U$ and $V$ such that $x\in U$, $y\in V$ and  $U \cap V = \varnothing$. Thus, $x \notin V$ and $y\notin U$, by Theorem \ref{T1th}, $X$ is also a $T_1$ space.
    
\end{proof}

\begin{theorem}\parencite{milewski1994topology}
    Each Hausdorff space is a $T_0$ space \label{hto}
\end{theorem}
\begin{proof}
    Suppose $X$ is a Hausdorff space. Then from Theorem \ref{hto1} and Theorem \ref{t1t0}, it is also a $T_0$ space.
\end{proof}


\begin{theorem} \parencite{Dugundji1966} The following three properties are equivalent

\begin{enumerate}
    \item $X$ is Hausdorff.
    \item Let $p \in X$. For each $q \ne p$, there is a nbd $U(p)$ such that $q \notin \overline{U(p)}$.
    \item For each $p \in X$, $\bigcap \{ \overline{U} \mid U \text{ is a nbd of } p \} = \{p\}$.
\end{enumerate}

\end{theorem}

\begin{theorem}
    Every subspace of a Hausdorff space is Hausdorff \label{subhaus}
\end{theorem}

\begin{theorem}
  If $(X,\tau)$ and $(Y,\tau')$ are homeomorphic  and $(X,\tau)$ is a $T_2$-space then so is  $(Y,\tau')$ 
\end{theorem}
\begin{proof}
    Let $f$ denote a homeomorphism, $f: X \to Y$ and $X$ be a $T_2$-space. Now $f$ is a continuous bijective map hence two distinct points $x_1,x_2$ of $X$ exist such that $f^{-1}(y_1)=x_1, f^{-1}(y_2)=x_2$. $(X, \tau)$ is a Hausdorff space, therefore, there are two open sets $U_1, U_2 \subset X$, such that
\[ x_1 \in U_1, \quad x_2 \in U_2, \quad U_1 \cap U_2 = \emptyset. \]
Since $f$ is bijective,
\[ f(U_1) \subset Y, \quad f(U_2) \subset Y \]
\[ f(U_1) \cap f(U_2) = \emptyset \]
Now, since $f^{-1}$ is continuous, the function $(f^{-1})^{-1} = f$ maps open sets into open sets. Hence, $f(U_1), f(U_2) \in T'$ are open sets.
\[ y_1 \in f(U_1), \quad y_2 \in f(U_2) \]
We conclude that $(Y, \tau')$ is a $T_2$-space.
If two spaces are homeomorphic and one of them is a $T_2$-space, then so is the other.
\end{proof}
 
\subsection{Regular Spaces}

\begin{definition} [Regular space] \label{regular}\parencite{milewski1994topology}
    A topological space $(X, \mathcal{T})$, is said to be regular if, given any closed subset $F \subset X$ and any point $x \in X$, such that $x \notin F$, there are open sets $U$ and $V$, such that
\[
F \subset U, \quad x \in V, \quad \text{and} \quad U \cap V = \emptyset 
\]
\end{definition}

\begin{definition} [$T_3$-space or Regular Hausdorff space]
A space is a $T_3$-space or Regular Hausdorff space if it is both a Hausdorff space and a regular space.
\end{definition}

Note that some authors switch the definition of "Regular" and "$T_3$". Some also defines them equivalently, such as \textcite{Dugundji1966}. Some  also defined $T_3$ by the theorems below. For this lecture notes, we use the definitions that was stated above.


\begin{remark}
A regular space need not be a $T_1$-space.    
\end{remark}

\begin{example}
Consider the topology $\tau = \{X, \emptyset, \{a\}, \{b, c\}\}$ on the set $X = \{a, b, c\}$. Observe that the closed subsets of $X$ are also $X$, $\emptyset$, $\{a\}$ and $\{b, c\}$ and that $(X, \tau)$ does satisfy Definition \ref{regular}. On the other hand, $(X, \tau)$ is not a $T_1$-space since there are single element sets, e.g. $\{b\}$, which are not closed.
\end{example}



\begin{theorem}
    A space is $T_3$ if and only if it is both regular and $T_0$
\end{theorem}
\begin{proof}
    Suppose a space is $T_3$ then by the definition, it is both regular and $T_2$. By Theorem \ref{hto}, it is also $T_0$. For the converse, suppose a space is both a regular and $T_0$. Now let $a,b \in X$ represents distinct points. Since the space is $T_0$, there is a neighborhood of at least one, which does not contain the other. Thus, let $a \notin U(b)$ then $a \notin \overline {U(b)}$. By regularity, there exist open sets $G$ and $V$ s.t $\overline {U(b)} \subset G$ and $a\in V$ where $G \cap V = \emptyset$. This implies that the space is also Hausdorff and hence also a $T_3$ space.
\end{proof}

\begin{theorem}
    A space is $T_3$ if and only if it is both regular and $T_1$
\end{theorem}

\begin{theorem} \parencite{Dugundji1966}
    The following three properties are equivalent
    \begin{enumerate}
        \item X is a $T_3$ space
         \item For each $x \in X$ and nbd $U$ of $x$, there exists a nbd $V$ of $x$ with $X \in V \subset \bar{V} \subset U$.
         \item For each $x \in X$ and closed $A$ not containing $x$, there is a nbd $V$ of $x$ with $\bar{V} \cap A = \emptyset$.
    \end{enumerate}
\end{theorem}

\begin{theorem}
    Every subspace of a regular space is regular \label{subreg}
\end{theorem}
\begin{proof}
 Let \( Y \) be a regular space and \( X \subset Y \) a subspace. Let \( B \subset X \) be closed in \( X \), and let \( x_0 \in X \setminus B \). Since \( B \) is closed in \( X \), there exists a closed set \( A \subset Y \) such that  
\( B = X \cap A \). Note that \( x_0 \notin A \) because \( x_0 \in X \setminus B \). By the regularity of \( Y \), there exist disjoint open sets \( U \) and \( V \) in \( Y \)  
such that \( x_0 \in U \) and \( A \subset V \). The sets \( U \cap X \) and \( V \cap X \) are open in \( X \) (by the definition of the subspace topology),  
\( x_0 \in U \cap X \), and  $B = X \cap A \subset V \cap X$. Since \( U \) and \( V \) are disjoint in \( Y \), their intersections with \( X \) are also disjoint in \( X \).

Thus, \( X \) is regular. 
\end{proof}

\begin{theorem}
    Any subspace of a $T_3$-space is a $T_3$-space
\end{theorem}
\begin{proof}
    A $T_3$-space is a regular Hausdorff space. By Theorem \ref{subreg} and theorem \ref{subhaus}, its subspace is also a regular Hausdorff space or a $T_3$-space.
\end{proof}

\subsection{Normal Spaces}
\begin{definition}[Normal Space] \parencite{milewski1994topology}
    A topological space $(X, T)$ is said to be normal if, given any two disjoint closed sets $F_1$ and $F_2$ in $X$, there are disjoint open sets $U$ and $V$, such that
\[ F_1 \subset U \quad \text{and} \quad F_2 \subset V \]
\end{definition}

\begin{definition}[$T_4$-space of Normal Hausdorff space]
A space is a $T_4$-space or Normal Hausdorff space if it is both a Hausdorff space and a normal space
\end{definition}

Similarly with regular spaces, some authors switch the definition of "Normal" and "$T_4$". Some also defines them equivalently, such as \textcite{Dugundji1966}. For this lecture notes, we use the definitions that are stated above.

\begin{example}
    Discrete spaces are $T_4$ spaces
\end{example}
\begin{proof}
For any two distinct points $x,y \in X$, the singleton sets $\{x\}$ and $\{y\}$ are open and disjoint. Thus, $X$ is Hausdorff. Now, Let $F_1$ and $F_2$ be disjoint closed sets in $X$. In the discrete topology, every set is open, so $F_1$ and $F_2$ themselves are open. Thus, we can take $U=F_1$ and $V=F_2$ as disjoint open sets containing $F_1$ and $F_2$ respectively. This shows that $X$ is normal. Since $X$ is both Hausdorff and normal, it is a $T_4$-space.    
\end{proof}

\begin{example}
    Any space $(X,\tau)$, containing more than one point with the indiscrete topology is Normal.
\end{example}

\begin{proof}
    The indiscrete topology consists of two sets $X$ and $\phi$.
\[ T = \{X, \phi\}. \]
Hence, the only closed sets are $X$ and $\phi$ because $X - \phi = X$ and $X - X = \phi$.
Thus, there are no non-empty disjoint closed subsets of $X$. The space is normal.
\end{proof}

\begin{theorem}
    Every $T_4$-spaces are $T_3$-spaces
\end{theorem}
\begin{proof}
Let $(X,T)$ denote a $T_4$-space. Hence, $(X,T)$ is normal and $T_1$. Suppose $F$ is a closed subset of $X$ and $a\in F$. Since $(X,T)$ is $T_1$, the singleton set $\{a\}$ is closed. Sets $F$ and $\{a\}$ are closed and disjoint. Since $(X,T)$ is normal, the open sets $U_1$ and $U_2$ exist, such that
$$\{a\} \subset U_1, \quad F\subset U_2, \quad U_1\cap U_2$$.
Therefore $(X,T)$ is regular and $T_3$
\end{proof}

\begin{theorem}
    The following four properties are equivalent
    \begin{enumerate}
        \item $X$ is $T_4$.
        \item  For each closed $A$ and open $U \supset A$ there is an open $V$ with $A \subset V \subset \bar{V} \subset U$.
        \item  For each pair of disjoint closed sets $A$, $B$, there is an open $U$ with $A \subset U$ and $\bar{U} \cap B = \emptyset$.
        \item Each pair of disjoint closed sets have nbds whose closures do not intersect.
    \end{enumerate}
\end{theorem}

%\begin{theorem}
 %    Normality is invariant under continuous closed surjections. 
%\end{theorem}

\begin{theorem}
    A closed subspace of a $T_4$ space is $T_4$.
\end{theorem}
\begin{proof} Let $X$ be a $T_4$ space and $Y$ be a closed subspace of $X$.
    Since every subspace of a $T_1$-space is $T_1$ and $X$ is $T_1$ also, $Y$ is a $T_1$-space. Since $Y$ is closed, a subset $F$ of $Y$ is closed in $Y$, if and only if $F$ is closed in $X$. Hence, if $F_1$ and $F_2$ are disjoint closed subsets of $Y$, they are also disjoint closed subsets of $X$.

Thus, the open sets $U_1$ and $U_2$ exist, such that
\[ F_1 \subset U_1, \quad F_2 \subset U_2 \text{ and } U_1 \cap U_2 = \emptyset. \]

Then
\[ F_1 \subset U_1 \cap Y, \quad F_2 \subset U_2 \cap Y, \]
and $U_1 \cap Y$ and $U_2 \cap Y$ are disjoint subsets of $Y$, open in $Y$. Since $(Y, T_Y)$ is $T_1$ and normal, it is $T_4$.
\end{proof}

\begin{remark}
    A subspace of a normal space need not be normal.
\end{remark}

\begin{definition}[ Completely normal space]
    A space $X$ is completely normal if every pair of sets $A, B$ satisfying $\bar{A} \cap B = A \cap \bar{B}=\emptyset$ can be separated. That is there exist disjoint open sets $U$ and $V$ such that $A \subseteq U$ and $B \subseteq V$.
\end{definition}

\begin{definition}[$T_5$-spaces or completely normal Hausdorff spaces] A space that is both Hausdorff and completely normal is a $T_5$ space.
\end{definition}


\subsection{Urysohn's Characterization of Normality}


\begin{theorem} [Urysohn Lemma]\parencite{Dugundji1966}
The following two properties are equivalent:
\begin{enumerate}
    \item $X$ is $T_4$.
    \item For each pair of disjoint closed sets, $A$, $B$ in $X$, there exists a continuous $f:X \to \mathbb{R}$, called a Urysohn function for $A$, $B$, such that:
    \begin{enumerate}
        \item $0 \leq f(x)\leq 1 \quad$ for all $x \in X$
        \item $f(a)=0$ \quad for all $a \in A$.
        \item $f(b)=0$ \quad for all $b \in B$.  
    \end{enumerate}
\end{enumerate}
\end{theorem}




\begin{corollary}\parencite{Dugundji1966} A necessary and sufficient condition for the existence of a Urysohn function satisfying $A=f^{-1}(0)$ is that $A$ be a $G_\delta$.
\end{corollary}

\begin{corollary}\parencite{Dugundji1966} A necessary and suffiecient condition that there be a Urysohn function $f$ with $A=f^{-1}(0), B=f^{-1}(1)$ is that both $A$ and $B$ be $G_\delta$.
\end{corollary}


\begin{definition}[$T_6$-spaces or perfectly normal Hausdorff spaces]
    A $T_4$ space in which each closed set is a $G_\delta$ is a $T_6$ space.
\end{definition}

\begin{theorem}
    Every $T_6$ space is a $T_5$ space.
\end{theorem}

\subsection{Tietze's Characterization of Normality}

\begin{theorem}[H. Tietze Theorem]\parencite{Dugundji1966} The following two properties are equivalent:
\begin{enumerate}
    \item $X$ is a $T_4$-space
    \item For every closed $A \subset X$, each continuous $f:A \to \mathbb{R}$ has a continuous $f:X\to \mathbb{R}$. Furthermore, if $|f(a)|<c$ on $A$, then $F$ can be chosen so that $|F(x)|<c$ on $X$.
\end{enumerate}
\end{theorem}

%\subsection{Covering Chracteization of Normality}

%\begin{theorem} The following two properties are equivalent:
%\begin{enumerate}
 %   \item Y is normal.
  %  \item If $\{U_\alpha \mid \alpha \in \mathcal{A}\}$ is any point-finite covering $Y$ by open sets, there exists a covering $\{V_\alpha \mid \alpha \in \mathcal{A}\}$ of $Y$ by open sets such that $\overline V_\alpha \subset U_\alpha$ for each $\alpha \in \mathcal{A}$, and $V_\alpha \neq \varnothing$ whenever $U_\alpha \neq \varnothing$.
%\end{enumerate}
% \end{theorem}

\subsection{Completely Regular Spaces}

\begin{definition}[Completely regular space]\parencite{Dugundji1966} A space is completely regular  if for each point $p\in X$ and closed $A$ not containing $p$, there is a continuous $\varphi:X \to [0,1]$ such that $\varphi(p)=1$ and $\varphi(a)=0$ for each $a\in A$
\end{definition}

\begin{theorem}
    Every completely regular space is regular \label{comples}
\end{theorem}
\begin{proof}
 Let $F$ represent a closed subset of $X$ and $a\in X$ a point which does not belong to $F$. By hypothesis, a continuous function
 $$ f: X \to [0,1]$$
 exists, such that $f(F) = \{1\}$ and $f(a) =0$. An interval $[0,1]$ is a Hausdorff space. Hence, two open disjoint subsets $U_1$ and $U_2$ of $[0,1]$ exists, such that
 $$ 0\in U_1 \quad \text{and} \quad  1 \in U_2.$$
 Since $f$ is continuous, $f^-1(U_1)$ and $f^-1(U_2)$ are open. These subsets are disjoint such that
 $$ a\in f^-1(U_1), \quad F\subset f^-1(U_2).$$
 Hence, $(X,T)$ is regular.
\end{proof}

\begin{definition}[Tychonoff space]
A completely regular Hausdorff space is a  Tychonoff space.
\end{definition}

\begin{theorem}
    \parencite{Dugundji1966} Every subspace of a Tychonoff space is Tychonoff. 
\end{theorem}
%\begin{proof}
 %   Suppose a space $X$ is a Tychonoff space. Then it is a completely regular Hausdorff Space.
%\end{proof}


%useful link https://www.emathzone.com/tutorials/general-topology/t0-space.html 


\newpage
\addcontentsline{toc}{chapter}{References}		
	\begingroup
            \printbibliography
        \endgroup
\end{document}